\section{引言：从预训练到 ChatGPT}

本节课由斯坦福大学博士生 Archit Sharma 主讲，主题是大语言模型的\textbf{后训练}（Post-training）。课程的核心问题是：我们如何从一个仅仅学习预测下一个 token 的预训练模型，一步步走到像 ChatGPT 这样能够理解用户意图、遵循指令、生成高质量回答的强大助手？

\begin{importantbox}{课程路线图}
本课覆盖四个核心主题，层层递进：
\begin{enumerate}
    \item \textbf{Zero-shot 与 Few-shot In-context Learning}：无需微调，通过提示词引导模型完成任务
    \item \textbf{指令微调（Instruction Fine-tuning / SFT）}：收集指令-回答对，训练模型遵循指令
    \item \textbf{基于人类偏好的优化（RLHF）}：引入人类偏好数据，通过强化学习对齐模型行为
    \item \textbf{直接偏好优化（DPO）}：绕过强化学习，直接用偏好数据优化语言模型
\end{enumerate}
\end{importantbox}

\subsection{预训练学到了什么？}

预训练的规模正在指数级增长。2022 年模型大约使用 1.4 万亿 token 进行训练，而到 2024 年，Llama 3 已经使用了约 \textbf{15 万亿 token}。模型消耗的计算量也从 $10^{24}$ FLOP 增长到超过 $10^{26}$ FLOP。

\begin{figure}[H]
\centering
\includegraphics[width=0.95\textwidth]{slides-images/slide-001.jpg}
\caption{预训练规模的指数级增长：计算量和数据量随时间急剧上升\protect\footnotemark}
\end{figure}
\footnotetext{Slides 第2页。}

尽管预训练的损失函数极其简单——仅仅是预测下一个 token，但模型在这个过程中学到的远不止语法和知识：

\begin{itemize}
    \item \textbf{世界知识}：``斯坦福大学位于加利福尼亚州''
    \item \textbf{数学能力}：理解方程和图形
    \item \textbf{代码编写}：Copilot 等代码补全工具
    \item \textbf{Agent 心智模型}：能够根据角色背景预测不同人物的行为
    \item \textbf{医学诊断}：提供初步的症状分析（但不应作为医疗建议）
\end{itemize}

\begin{figure}[H]
\centering
\includegraphics[width=0.95\textwidth]{slides-images/slide-004.jpg}
\caption{预训练中学到的 Agent 心智模型：改变角色描述后，模型的预测也随之改变\protect\footnotemark}
\end{figure}
\footnotetext{Slides 第5页。}

\begin{knowledgebox}{预训练不只是记忆}
一个有趣的实验：当文本描述``Pat 是一名物理学家''时，模型预测保龄球和树叶会同时落地（正确的物理知识）；而当描述改为``Pat 从未见过这个实验''时，模型预测保龄球先落地（符合直觉但物理上不正确）。这说明模型不仅记住了事实，还学会了根据\textbf{角色背景}进行推理——这种能力远超简单的文本预测。
\end{knowledgebox}

\begin{figure}[H]
\centering
\includegraphics[width=0.95\textwidth]{slides-images/slide-005.jpg}
\caption{预训练中涌现的多种能力：数学、代码、医学等\protect\footnotemark}
\end{figure}
\footnotetext{Slides 第6页。}

语言模型正在发展成为\textbf{通用多任务助手}。从简单的文本预测出发，它们已经能够处理代码、数学、创意写作、医疗咨询等多种任务。今天的课程将解释如何通过后训练技术进一步释放和对齐这些能力。

\subsection{本章小结}

预训练通过海量数据上的 next-token prediction 为模型注入了丰富的知识和推理能力。然而，预训练模型本质上是一个文本补全器，并非一个助手。后训练的目标就是将这个强大但``不听话''的补全器，转变为一个能理解用户意图、遵循指令、安全可靠的 AI 助手。

%% ========== Part 2: Zero-shot & Few-shot ==========

